% generated by tools/cp.py setup -- do not edit
\section{Template}
\begin{libentry}{Full template}{Ctrl+Alt+T, T}{cptemplate}{empty file}{--}{template/template.cpp}
the whole template, for files CPH didn\textquotesingle{}t create
\end{libentry}
\section{Basics: input, search, sorting, prefix sums, two pointers}
\begin{libentry}{2D prefix sums}{Ctrl+Alt+T, Shift+A}{cpprefix2d}{inside solve()}{O(nm) build, O(1) per query}{library/basics/prefix\_sum\_2d.cpp}
needs n, m, a (n x m). rect(r1, c1, r2, c2) = sum of rows [r1, r2) x cols [c1, c2).
\begin{lstlisting}
ll s = rect(r1, c1, r2, c2);  // rows [r1, r2), cols [c1, c2)
\end{lstlisting}
\end{libentry}
\begin{libentry}{All permutations (next\_permutation)}{Ctrl+Alt+T, 9}{cpperm}{inside solve()}{O(n! n)}{library/basics/permutations.cpp}
needs n (n \textless{}= \textasciitilde{}10). Visits every order of 0..n-1 in lexicographic order.
\begin{lstlisting}
// on a sorted vector v directly: do { ... } while (next_permutation(v.begin(), v.end()));  (skips duplicates)
\end{lstlisting}
\end{libentry}
\begin{libentry}{All subsets (bitmask loop)}{Ctrl+Alt+T, 8}{cpsubsets}{inside solve()}{O(2\textasciicircum{}n n)}{library/basics/subsets.cpp}
needs n, a (n \textless{}= \textasciitilde{}20). Visits every subset; bit i of mask set = item i taken.
\begin{lstlisting}
// for all submasks of a mask: for (int s = mask; s; s = (s - 1) & mask)
\end{lstlisting}
\end{libentry}
\begin{libentry}{Argsort + custom comparator}{Ctrl+Alt+T, Shift+J}{cpargsort}{above solve()}{O(n log n)}{library/basics/argsort.cpp}
argsort(a): indices sorted by a[i] (stable). Custom orders: sort with a lambda (see usage).
\begin{lstlisting}
vector<int> idx = argsort(a);  // a[idx[0]] is the smallest
sort(v.begin(), v.end(), [](const auto& x, const auto& y) { return make_pair(x.first, -x.second) < make_pair(y.first, -y.second); });
\end{lstlisting}
\end{libentry}
\begin{libentry}{Binary search on a real answer}{Ctrl+Alt+T, Shift+F}{cpbsreal}{inside solve()}{100 calls to ok}{library/basics/binary\_search\_real.cpp}
smallest real x in [lo, hi] with ok(x); ok must be monotone. 100 halvings reach full double precision.
\begin{lstlisting}
// fill in ok(x); lo ends within 1e-15 of the boundary (relative to hi - lo)
\end{lstlisting}
\end{libentry}
\begin{libentry}{Binary search on the answer}{Ctrl+Alt+T, A}{cpbsearch}{inside solve()}{O(log(hi - lo)) calls to ok}{library/basics/binary\_search.cpp}
smallest x in [lo, hi] with ok(x); ok must be monotone (false...false true...true).
\begin{lstlisting}
// fill in ok(x); lo ends as the smallest x with ok(x)
\end{lstlisting}
\end{libentry}
\begin{libentry}{Bit helpers}{Ctrl+Alt+T, Y}{cpbits}{above solve()}{O(1)}{library/basics/bits.cpp}
popcnt, highest\_bit (floor log2), lowest\_bit, has\_bit for 64-bit values; submask loop in usage.
\begin{lstlisting}
popcnt(x); highest_bit(x); lowest_bit(x); has_bit(x, i);
for (ll s = mask; s; s = (s - 1) & mask) { /* every non-empty submask */ }
\end{lstlisting}
\end{libentry}
\begin{libentry}{Coordinate compression}{Ctrl+Alt+T, Shift+C}{cpcompress}{above solve()}{O(n log n)}{library/basics/compress.cpp}
auto vals = compress(a); a[i] becomes its rank 0..k-1 and vals[rank] is the original value.
\begin{lstlisting}
vector<int> vals = compress(a);  // a[i] -> rank, vals[rank] -> value
\end{lstlisting}
\end{libentry}
\begin{libentry}{Difference array (offline range add, 1D)}{Ctrl+Alt+T, 3}{cpdiff}{inside solve()}{O(1) per update, O(n) build}{library/basics/difference\_array.cpp}
needs n. range\_add(l, r, v) adds v to a[l, r) in O(1); build() returns the final array.
\begin{lstlisting}
range_add(l, r, v);            // any number of times
vector<ll> a = build();        // once, at the end
\end{lstlisting}
\end{libentry}
\begin{libentry}{Difference array (offline rectangle add, 2D)}{Ctrl+Alt+T, 4}{cpdiff2d}{inside solve()}{O(1) per update, O(nm) build}{library/basics/difference\_array\_2d.cpp}
needs n, m. rect\_add(r1, c1, r2, c2, v) adds v to rows [r1, r2) x cols [c1, c2); build() returns the grid.
\begin{lstlisting}
rect_add(r1, c1, r2, c2, v);   // any number of times
auto g = build();              // g[r][c] = total added to cell (r, c)
\end{lstlisting}
\end{libentry}
\begin{libentry}{Gray code}{Ctrl+Alt+U, 0}{cpgray}{above solve()}{O(1) / O(log) per value}{library/basics/gray\_code.cpp}
gray(i) = i \textasciicircum{} (i \textgreater{}\textgreater{} 1): consecutive values differ in exactly one bit; gray\_inverse(g) gives i back. To visit all 2\textasciicircum{}n subsets changing one element at a time, step k flips bit \_\_builtin\_ctzll(k).
\begin{lstlisting}
for (ll k = 1, mask = 0; k < (1LL << n); k++) { int b = __builtin_ctzll(k); mask ^= 1LL << b; /* add or remove b */ }
ll g = gray(i), back = gray_inverse(g);
\end{lstlisting}
\end{libentry}
\begin{libentry}{Intervals: merge + max disjoint (greedy)}{Ctrl+Alt+T, Shift+N}{cpintervals}{above solve()}{O(n log n)}{library/basics/intervals.cpp}
merge\_intervals(v): union of closed [l, r] intervals; max\_disjoint(v): most pairwise non-overlapping (earliest end first).
\begin{lstlisting}
vector<pair<ll, ll>> segs = merge_intervals(v);
int k = max_disjoint(v);  // touching intervals ([1, 2] and [2, 3]) count as overlapping
\end{lstlisting}
\end{libentry}
\begin{libentry}{Meet in the middle (subset sums)}{Ctrl+Alt+T, Shift+U}{cpmitm}{above solve()}{O(2\textasciicircum{}(n/2) n)}{library/basics/meet\_in\_middle.cpp}
count\_subsets\_leq(a, X): number of subsets with sum \textless{}= X for n \textless{}= \textasciitilde{}40, splitting the items in two halves.
\begin{lstlisting}
ll ways = count_subsets_leq(a, X);
// exact sum X: count upper_bound - lower_bound of X - x in the sorted right half
\end{lstlisting}
\end{libentry}
\begin{libentry}{Monotonic stack (next greater / previous smaller)}{Ctrl+Alt+T, 7}{cpmonostack}{above solve()}{O(n)}{library/basics/monotonic\_stack.cpp}
next\_greater(a)[i] = first j \textgreater{} i with a[j] \textgreater{} a[i] (n if none); prev\_smaller(a)[i] = last j \textless{} i with a[j] \textless{} a[i] (-1).
\begin{lstlisting}
auto nxt = next_greater(a); auto prv = prev_smaller(a);
// largest rectangle in a histogram: width of bar i = next smaller - prev smaller - 1
\end{lstlisting}
\end{libentry}
\begin{libentry}{Prefix sums (1D)}{Ctrl+Alt+T, 1}{cpprefix}{inside solve()}{O(n) build, O(1) per query}{library/basics/prefix\_sum.cpp}
needs n, a. P[i] = a[0] + ... + a[i-1]; range\_sum(l, r) = a[l] + ... + a[r-1].
\begin{lstlisting}
ll s = range_sum(l, r);  // [l, r)
// subarrays with sum k: count pairs i < j with P[j] - P[i] == k using a map of seen P values
\end{lstlisting}
\end{libentry}
\begin{libentry}{Radix sort and counting sort}{Ctrl+Alt+U, Shift+R}{cpradix}{above solve()}{O(n) per pass: O(4 (n + 65536)) for radix\_sort, O(n + K) for counting\_order}{library/basics/radix\_sort.cpp}
radix\_sort(a): sorts 64-bit integers (negatives too) in 4 linear passes of 16 bits, often 2-3x faster than std::sort for n \textgreater{}= 1e6; counting\_order(keys, K): stable order of indices by small keys in [0, K).
\begin{lstlisting}
radix_sort(a);   // vector<ll>
vector<int> ord = counting_order(keys, K);   // stable: equal keys keep their order
\end{lstlisting}
\end{libentry}
\begin{libentry}{Read graph (n, m, m edges)}{Ctrl+Alt+T, Shift+G}{cpgraph}{inside solve()}{O(n + m)}{library/basics/read\_graph.cpp}
reads n, m and m 1-indexed edges into a 0-indexed undirected adjacency list.
\begin{lstlisting}
// then: dfs / bfs / ... over adj
\end{lstlisting}
\end{libentry}
\begin{libentry}{Read vector}{Ctrl+Alt+T, V}{cpread}{inside solve()}{O(n)}{library/basics/read\_vector.cpp}
needs n. Reads n numbers into a.
\begin{lstlisting}
// with the template: vi a(n); cin >> a;  does the same
\end{lstlisting}
\end{libentry}
\begin{libentry}{Read weighted graph (n, m, edges u v w)}{Ctrl+Alt+T, Shift+E}{cpgraphw}{inside solve()}{O(n + m)}{library/basics/read\_weighted\_graph.cpp}
reads n, m and m 1-indexed weighted edges into adj[u] = \{v, w\}, 0-indexed, undirected.
\begin{lstlisting}
// then: auto [dist, par] = dijkstra(adj, 0);
\end{lstlisting}
\end{libentry}
\begin{libentry}{Recursive complete search (backtracking)}{Ctrl+Alt+T, Shift+Y}{cpbacktrack}{inside solve()}{O(2\textasciicircum{}n) calls without pruning}{library/basics/backtracking.cpp}
needs n. Builds every subsequence of 0..n-1 in cur (skip / take); prune at the top of search().
\begin{lstlisting}
// permutations instead: loop over unused items, mark used, recurse, unmark
\end{lstlisting}
\end{libentry}
\begin{libentry}{Sliding window min / max (monotonic deque)}{Ctrl+Alt+T, 6}{cpwindow}{above solve()}{O(n)}{library/basics/sliding\_window.cpp}
window\_min(a, k)[i] = min of a[i, i + k); pass want\_max = true for the max. Monotonic deque.
\begin{lstlisting}
vector<ll> mins = window_min(a, k);
vector<ll> maxs = window_min(a, k, true);
\end{lstlisting}
\end{libentry}
\begin{libentry}{Ternary search (unimodal max)}{Ctrl+Alt+T, Q}{cpternary}{above solve()}{O(log range) calls to f}{library/basics/ternary\_search.cpp}
argmax\_int(lo, hi, f) / argmax\_real(lo, hi, f) for f that rises then falls. For a minimum pass -f.
\begin{lstlisting}
ll x = argmax_int(0, 1e9, [&](ll x) { return profit(x); });
ld y = argmax_real(0, 1e4, [&](ld x) { return -cost(x); });  // minimum of cost
\end{lstlisting}
\end{libentry}
\begin{libentry}{Two pointers (sliding window)}{Ctrl+Alt+T, 5}{cptwoptr}{inside solve()}{O(n)}{library/basics/two\_pointers.cpp}
needs n, a (values \textgreater{}= 0). Longest window with sum \textless{}= S; grow r, shrink l while the window is invalid.
\begin{lstlisting}
// count windows instead: add (r - l + 1) to an answer inside the loop
\end{lstlisting}
\end{libentry}
\section{Graphs}
\begin{libentry}{0-1 BFS}{Ctrl+Alt+T, 0}{cp01bfs}{above solve()}{O(n + m)}{library/graph/zero\_one\_bfs.cpp}
vector\textless{}ll\textgreater{} dist = zero\_one\_bfs(adj, src); adj[u] = \{v, w\} with w in \{0, 1\}. Unreachable = LLONG\_MAX.
\begin{lstlisting}
vector<ll> dist = zero_one_bfs(adj, src);  // adj[u] = {v, w}, w in {0, 1}
\end{lstlisting}
\end{libentry}
\begin{libentry}{2-edge-connected components}{Ctrl+Alt+Y, B}{cp2ecc}{above solve()}{O(n + m)}{library/graph/two\_edge\_cc.cpp}
auto [k, comp] = two\_edge\_cc(n, edges): vertices stay connected after deleting any single edge iff same comp. Bridges join comps.
\begin{lstlisting}
auto [k, comp] = two_edge_cc(n, edges);
// bridge tree: for each edge (u, v) with comp[u] != comp[v], connect comp[u] - comp[v]
\end{lstlisting}
\end{libentry}
\begin{libentry}{2-SAT}{Ctrl+Alt+Y, 2}{cp2sat}{above solve()}{O(n + clauses)}{library/graph/two\_sat.cpp}
TwoSat ts(n); ts.either(a, va, b, vb) adds the clause (x\_a == va OR x\_b == vb); ts.solve() -\textgreater{} values or empty if impossible.
\begin{lstlisting}
TwoSat ts(n);
ts.either(a, true, b, false); ts.implies(a, true, c, true); ts.set(d, false);
vector<int> x = ts.solve();  // empty: unsatisfiable
\end{lstlisting}
\end{libentry}
\begin{libentry}{Bellman-Ford (negative edges)}{Ctrl+Alt+T, Shift+W}{cpbellman}{above solve()}{O(nm)}{library/graph/bellman\_ford.cpp}
bellman\_ford(n, edges, src), edges = \{u, v, w\}. BF\_INF = unreachable, -BF\_INF = hit by a negative cycle.
\begin{lstlisting}
vector<array<ll, 3>> edges;  // {u, v, w}
vector<ll> dist = bellman_ford(n, edges, src);  // -BF_INF: negative cycle reaches it
\end{lstlisting}
\end{libentry}
\begin{libentry}{BFS (shortest path, unweighted)}{Ctrl+Alt+T, B}{cpbfs}{inside solve()}{O(n + m)}{library/graph/bfs.cpp}
needs n, adj, src. dist[v] = -1 if unreachable; get\_path(t) rebuilds the route.
\begin{lstlisting}
// after it: dist[v] (-1 = unreachable), par[v]
vector<int> path = get_path(t);  // src ... t
\end{lstlisting}
\end{libentry}
\begin{libentry}{Biconnected components (blocks)}{Ctrl+Alt+Y, Shift+B}{cpbcc}{above solve()}{O(n + m)}{library/graph/biconnected\_components.cpp}
biconnected\_components(n, edges): every block (maximal subgraph with no cut vertex) as a vertex list. Cut vertices appear in several.
\begin{lstlisting}
auto blocks = biconnected_components(n, edges);
// block-cut tree: node per block + node per cut vertex, edge when the cut vertex is in the block
\end{lstlisting}
\end{libentry}
\begin{libentry}{Bidirectional BFS (meet in the middle on shortest paths)}{Ctrl+Alt+U, 2}{cpbibfs}{above solve()}{O(b\textasciicircum{}(d/2)) states for branching factor b and distance d, instead of O(b\textasciicircum{}d)}{library/graph/bidirectional\_bfs.cpp}
bidirectional\_bfs(s, t, neighbours): unweighted distance from s to t (-1 if unreachable) exploring from both ends; neighbours(state) returns a vector of states. Works on implicit graphs (puzzles, strings) with any hashable State.
\begin{lstlisting}
ll d = bidirectional_bfs<ll>(start, goal, [&](ll s) { vector<ll> nb; /* moves */ return nb; });
\end{lstlisting}
\end{libentry}
\begin{libentry}{Bipartite check (2-coloring)}{Ctrl+Alt+T, 2}{cpbipartite}{above solve()}{O(n + m)}{library/graph/bipartite.cpp}
vector\textless{}int\textgreater{} color = bipartite(adj); colors are 0/1, empty if an odd cycle exists.
\begin{lstlisting}
vector<int> color = bipartite(adj);
if (color.empty()) { /* odd cycle: not bipartite */ }
\end{lstlisting}
\end{libentry}
\begin{libentry}{Bridges + articulation points}{Ctrl+Alt+T, Shift+B}{cpbridges}{above solve()}{O(n + m)}{library/graph/bridges.cpp}
auto [bridge\_ids, cut\_vertices] = bridges\_and\_cuts(n, edges); undirected, multi-edges OK.
\begin{lstlisting}
vector<pair<int, int>> edges;  // undirected {u, v}
auto [bridge_ids, cut_vertices] = bridges_and_cuts(n, edges);
\end{lstlisting}
\end{libentry}
\begin{libentry}{Cycle finding in O(1) memory (Floyd / Brent)}{Ctrl+Alt+U, F}{cpcyclefind}{above solve()}{O(mu + lambda) evaluations of f, O(1) memory}{library/graph/cycle\_floyd.cpp}
find\_rho(x0, f): for the sequence x0, f(x0), f(f(x0)), ... returns \{mu, lambda\}: the first index of the cycle and its length. Brent\textquotesingle{}s variant: O(mu + lambda) calls of f, O(1) memory (use when states are huge).
\begin{lstlisting}
auto [mu, lam] = find_rho<ll>(seed, [](ll x) { return (x * x + 1) % 1'000'003; });
\end{lstlisting}
\end{libentry}
\begin{libentry}{DFS (recursive, all components)}{Ctrl+Alt+T, D}{cpdfs}{inside solve()}{O(n + m)}{library/graph/dfs.cpp}
needs n, adj. Labels every vertex with its component id.
\begin{lstlisting}
// after it: comp[v] = component of v, ncomp = number of components
\end{lstlisting}
\end{libentry}
\begin{libentry}{Dijkstra}{Ctrl+Alt+T, J}{cpdijkstra}{above solve()}{O((n + m) log n)}{library/graph/dijkstra.cpp}
auto [dist, par] = dijkstra(adj, src); adj[u] = \{v, w\}, w \textgreater{}= 0. Unreachable = LLONG\_MAX.
\begin{lstlisting}
vector<vector<pair<int, ll>>> adj(n);  // adj[u].push_back({v, w})
auto [dist, par] = dijkstra(adj, src);
\end{lstlisting}
\end{libentry}
\begin{libentry}{DSU (union-find)}{Ctrl+Alt+T, U}{cpdsu}{above solve()}{\textasciitilde{}O(1) amortized per call}{library/graph/dsu.cpp}
DSU d(n); d.unite(a, b) -\textgreater{} false if already joined; d.same, d.size, d.comps.
\begin{lstlisting}
DSU d(n);
if (d.unite(a, b)) { /* merged two components */ }
d.same(a, b); d.size(a); d.comps;
\end{lstlisting}
\end{libentry}
\begin{libentry}{DSU with rollback}{Ctrl+Alt+Y, R}{cpdsurb}{above solve()}{O(log n) per find / unite, O(1) per undone unite}{library/graph/dsu\_rollback.cpp}
DSURollback d(n); d.unite(a, b); int t = d.snapshot(); ... d.rollback(t) undoes every unite after t. No path compression.
\begin{lstlisting}
DSURollback d(n);
int t = d.snapshot(); d.unite(a, b); /* ... */ d.rollback(t);
\end{lstlisting}
\end{libentry}
\begin{libentry}{Eulerian path / circuit (Hierholzer)}{Ctrl+Alt+Y, E}{cpeuler}{above solve()}{O(n + m)}{library/graph/euler\_path.cpp}
euler\_path(n, edges, directed, src = -1): vertex sequence using every edge exactly once, or empty if impossible.
\begin{lstlisting}
vector<int> path = euler_path(n, edges, false);       // path.size() == m + 1
vector<int> tour = euler_path(n, edges, false, 0);    // circuit from 0 (only when every degree is even)
\end{lstlisting}
\end{libentry}
\begin{libentry}{Find a directed cycle}{Ctrl+Alt+T, /}{cpcycle}{above solve()}{O(n + m)}{library/graph/find\_cycle.cpp}
find\_cycle(adj): vertices of some directed cycle in order (the last one has an edge back to the first); empty if acyclic.
\begin{lstlisting}
vector<int> cyc = find_cycle(adj);
// undirected graphs: a DSU finds the first edge whose ends are already joined
\end{lstlisting}
\end{libentry}
\begin{libentry}{Flood fill (grid regions)}{Ctrl+Alt+T, F}{cpflood}{inside solve()}{O(nm)}{library/graph/flood\_fill.cpp}
needs n, m, grid (vector\textless{}string\textgreater{}). Iterative, so big grids can\textquotesingle{}t overflow the stack.
\begin{lstlisting}
// after it: region[r][c] = region id (-1 = not open), region_size[id]
// change open(r, c) to choose which cells get filled
\end{lstlisting}
\end{libentry}
\begin{libentry}{Floyd-Warshall (all pairs)}{Ctrl+Alt+T, W}{cpfloyd}{above solve()}{O(n\textasciicircum{}3)}{library/graph/floyd\_warshall.cpp}
floyd\_warshall(d) in place, O(n\textasciicircum{}3). d[i][i] \textless{} 0 afterwards \textless{}=\textgreater{} i is on a negative cycle.
\begin{lstlisting}
vector<vector<ll>> d(n, vector<ll>(n, FW_INF));
for (int i = 0; i < n; i++) d[i][i] = 0;  // then d[u][v] = min(d[u][v], w) per edge
floyd_warshall(d);
\end{lstlisting}
\end{libentry}
\begin{libentry}{Functional graph (successor graph) cycles}{Ctrl+Alt+T, ,}{cpfuncgraph}{above solve()}{O(n)}{library/graph/functional\_graph.cpp}
FunctionalGraph fg(f), f[v] = the one successor of v: fg.cyc[v] (cycle id or -1), fg.cyc\_len[id], fg.dist[v] to its cycle, fg.entry[v].
\begin{lstlisting}
FunctionalGraph fg(f);
int steps_until_repeat = fg.dist[v] + fg.cyc_len[fg.cyc[fg.entry[v]]];
\end{lstlisting}
\end{libentry}
\begin{libentry}{Grid BFS (shortest path on a grid)}{Ctrl+Alt+T, G}{cpgridbfs}{inside solve()}{O(nm)}{library/graph/grid\_bfs.cpp}
needs n, m, grid. Starts from every \textquotesingle{}S\textquotesingle{} (multi-source), \textquotesingle{}\#\textquotesingle{} is a wall.
\begin{lstlisting}
// after it: dist[r][c] = steps from the nearest 'S' (-1 = unreachable)
\end{lstlisting}
\end{libentry}
\begin{libentry}{K-th successor (binary lifting / jumping)}{Ctrl+Alt+T, .}{cpjump}{above solve()}{O(n log K) build, O(log K) per query}{library/graph/kth\_successor.cpp}
Jump J(f, max\_k); J.go(v, k) = f applied k times to v, for k \textless{}= max\_k. Works for any functional graph.
\begin{lstlisting}
Jump J(f, 1e9);
int where = J.go(v, k);
\end{lstlisting}
\end{libentry}
\begin{libentry}{Kruskal MST}{Ctrl+Alt+T, K}{cpkruskal}{above solve()}{O(m log m)}{library/graph/kruskal.cpp}
auto [cost, used] = kruskal(n, edges); edges = \{u, v, w\}. Spanning forest if disconnected.
\begin{lstlisting}
vector<array<ll, 3>> edges;  // {u, v, w}
auto [cost, used] = kruskal(n, edges);
\end{lstlisting}
\end{libentry}
\begin{libentry}{Offline dynamic connectivity (edge deletions)}{Ctrl+Alt+Y, Shift+R}{cpofflinecc}{above solve()}{O(q log q log n)}{library/graph/offline\_connectivity.cpp}
offline\_components(n, ops): ops are \{0, u, v\} add edge, \{1, u, v\} remove it, \{2, 0, 0\} query; returns \#components per query. Needs cpdsurb.
\begin{lstlisting}
vector<array<int, 3>> ops;  // {type, u, v}
vector<int> answers = offline_components(n, ops);
\end{lstlisting}
\end{libentry}
\begin{libentry}{Prim MST (dense graphs)}{Ctrl+Alt+T, Shift+5}{cpprim}{above solve()}{O(n\textasciicircum{}2)}{library/graph/prim.cpp}
prim\_dense(w): MST weight for an n x n weight matrix (complete graphs, e.g. points in the plane). O(n\textasciicircum{}2).
\begin{lstlisting}
vector<vector<ll>> w(n, vector<ll>(n));  // w[i][j] = cost of edge i-j
ll cost = prim_dense(w);
\end{lstlisting}
\end{libentry}
\begin{libentry}{SCC (Tarjan)}{Ctrl+Alt+T, C}{cpscc}{above solve()}{O(n + m)}{library/graph/scc.cpp}
auto [k, comp] = scc(adj). Ids are reverse-topological: edge u-\textgreater{}v across SCCs has comp[u] \textgreater{} comp[v].
\begin{lstlisting}
auto [k, comp] = scc(adj);  // k components, comp[v] in [0, k)
\end{lstlisting}
\end{libentry}
\begin{libentry}{System of difference constraints}{Ctrl+Alt+U, Shift+I}{cpdiffcons}{above solve()}{O(n m) (Bellman-Ford from a virtual source)}{library/graph/difference\_constraints.cpp}
DiffConstraints dc(n); dc.add(u, v, w) means x\_v - x\_u \textless{}= w; dc.solve() -\textgreater{} a solution vector (all x \textless{}= 0, maximal), or empty if infeasible (a negative cycle).
\begin{lstlisting}
DiffConstraints dc(n); dc.add(a, b, 3);   // x_b <= x_a + 3
vector<ll> x = dc.solve(); if (x.empty()) { /* impossible */ }
\end{lstlisting}
\end{libentry}
\begin{libentry}{Topological sort (Kahn)}{Ctrl+Alt+T, O}{cptopo}{above solve()}{O(n + m)}{library/graph/topo\_sort.cpp}
vector\textless{}int\textgreater{} order = topo\_sort(adj); empty if the graph has a cycle.
\begin{lstlisting}
vector<int> order = topo_sort(adj);
if (order.empty()) { /* the graph has a cycle */ }
\end{lstlisting}
\end{libentry}
\section{Trees}
\begin{libentry}{Centroid decomposition}{Ctrl+Alt+Y, C}{cpcentroid}{above solve()}{O(n log n)}{library/tree/centroid.cpp}
CentroidDecomp cd(adj): cd.cpar[v] = parent in the centroid tree (-1 at the top), cd.clevel[v] = its depth (\textless{}= log2 n).
\begin{lstlisting}
CentroidDecomp cd(adj);
for (int c = v; c != -1; c = cd.cpar[c]) { /* every path through v passes a centroid ancestor c: use dist(v, c) */ }
\end{lstlisting}
\end{libentry}
\begin{libentry}{DSU on tree (Sack): subtree statistics in one shared array}{Ctrl+Alt+U, Shift+T}{cpsack}{inside solve()}{O(n log n) calls of add}{library/tree/sack.cpp}
needs n, adj, root, color (values in [0, n)). Keeps the heavy child\textquotesingle{}s counts and re-adds the light subtrees: best[v] = the largest count of one color in v\textquotesingle{}s subtree. Edit add / clear / the answer line.
\begin{lstlisting}
// add(v) puts one vertex into cnt[] and the statistic; clear(v) empties the subtree of v again
\end{lstlisting}
\end{libentry}
\begin{libentry}{Heavy-light decomposition}{Ctrl+Alt+Y, H}{cphld}{above solve()}{O(n) build, O(log n) ranges per path}{library/tree/hld.cpp}
HLD h(adj, root); h.for\_path(a, b, f) calls f(l, r) on position ranges [l, r) covering the path; pair with any [l, r) segment tree.
\begin{lstlisting}
HLD h(adj); SegTree st(n);                      // st.set(h.pos[v], value of v)
ll s = 0; h.for_path(a, b, [&](int l, int r) { s += st.query(l, r); });
// subtree of v = [h.pos[v], h.pos[v] + h.sub[v]); h.lca(a, b)
\end{lstlisting}
\end{libentry}
\begin{libentry}{Jump pointers (level ancestors with one extra pointer per vertex)}{Ctrl+Alt+U, J}{cpjumpptr}{above solve()}{O(1) per added leaf, O(log n) per query}{library/tree/jump\_pointers.cpp}
JumpPointers jp(n); jp.add\_root(r); jp.add\_leaf(v, p) as vertices arrive (parent first); jp.kth(v, k), jp.ancestor\_at\_depth(v, d), jp.lca(a, b). Works online while the tree grows.
\begin{lstlisting}
JumpPointers jp(n); jp.add_root(0);
for (int v : order) if (v != 0) jp.add_leaf(v, par[v]);  // any order with parents before children
int a = jp.kth(v, k);  // -1 if v has fewer than k ancestors
\end{lstlisting}
\end{libentry}
\begin{libentry}{Kruskal reconstruction tree}{Ctrl+Alt+Y, K}{cpkrt}{above solve()}{O(m log m)}{library/tree/kruskal\_reconstruction.cpp}
KRT t(n, edges \{u, v, w\}): leaves 0..n-1, each merge makes node x with val[x] = w. Max edge on the MST path a-b = val[lca(a, b)].
\begin{lstlisting}
KRT t(n, edges); auto g = t.tree();  // then LCA l(g, g.size() - 1) if the graph is connected
// vertices reachable from v using edges <= x = leaves under v's highest ancestor with val <= x
\end{lstlisting}
\end{libentry}
\begin{libentry}{LCA (binary lifting)}{Ctrl+Alt+T, L}{cplca}{above solve()}{O(n log n) build, O(log n) per query}{library/tree/lca.cpp}
LCA t(adj, root); t.lca(a, b), t.dist(a, b), t.kth\_ancestor(v, k), t.depth[v].
\begin{lstlisting}
LCA t(adj, root);
t.lca(a, b); t.dist(a, b); t.kth_ancestor(v, k); t.depth[v];
\end{lstlisting}
\end{libentry}
\begin{libentry}{Long-path decomposition: depth-indexed tree DP in O(n)}{Ctrl+Alt+U, Shift+N}{cplongpath}{inside solve()}{O(n)}{library/tree/long\_path.cpp}
needs n, adj, root. f\_v[d] = \# vertices at distance d below v lives at buf[pos[v] + d], shared with the long child (shifted by one); dom[v] = the smallest d maximising f\_v[d]. Edit the merge for other depth DPs.
\begin{lstlisting}
// the array of v has h[v] cells: buf[pos[v]] .. buf[pos[v] + h[v] - 1]
\end{lstlisting}
\end{libentry}
\begin{libentry}{Offline LCA (Tarjan)}{Ctrl+Alt+U, Shift+L}{cpofflinelca}{above solve()}{O((n + q) alpha(n))}{library/tree/offline\_lca.cpp}
tarjan\_lca(adj, root, qs): the LCA of every query pair, answered in one DFS with a DSU. Faster than binary lifting when all queries are known in advance and memory is tight.
\begin{lstlisting}
vector<int> ans = tarjan_lca(adj, 0, qs);   // qs: vector<pair<int, int>>
\end{lstlisting}
\end{libentry}
\begin{libentry}{Prufer code (tree \textless{}-\textgreater{} sequence)}{Ctrl+Alt+U, U}{cpprufer}{above solve()}{O(n) (linear pointer technique)}{library/tree/prufer.cpp}
prufer\_encode(n, edges): the Prufer sequence (length n-2) of a labelled tree on 0..n-1; prufer\_decode(seq): the tree back. Bijection behind Cayley\textquotesingle{}s formula n\textasciicircum{}(n-2); vertex v appears deg(v)-1 times.
\begin{lstlisting}
vector<int> seq = prufer_encode(n, edges);
vector<pair<int, int>> tree = prufer_decode(seq);
\end{lstlisting}
\end{libentry}
\begin{libentry}{Rerooting (DP for every root)}{Ctrl+Alt+T, \textquotesingle{}}{cpreroot}{above solve()}{O(n) merges and lifts}{library/tree/rerooting.cpp}
reroot\textless{}T\textgreater{}(adj, id, merge, lift): res[v] = merge over v\textquotesingle{}s neighbours w of lift(value of w\textquotesingle{}s side, w). merge associative, id its identity.
\begin{lstlisting}
// farthest vertex from every v:  auto far = reroot<ll>(adj, 0, [](ll a, ll b) { return max(a, b); }, [](ll x, int) { return x + 1; });
// sum of distances: T = pair<ll, ll> {sum, count}, lift(p, w) = {p.first + p.second + 1, p.second + 1}, merge = add
\end{lstlisting}
\end{libentry}
\begin{libentry}{Small-to-large merging}{Ctrl+Alt+Y, M}{cpsmall2large}{inside solve()}{O(n log\textasciicircum{}2 n)}{library/tree/small\_to\_large.cpp}
needs n, adj, root, color. distinct[v] = \# distinct colors in v\textquotesingle{}s subtree; always pour the smaller set into the larger.
\begin{lstlisting}
// swap set<int> for map<int, int> to keep counts (e.g. most frequent color per subtree)
\end{lstlisting}
\end{libentry}
\begin{libentry}{Tree DFS (parent, depth, subtree size, Euler tour)}{Ctrl+Alt+T, E}{cptree}{inside solve()}{O(n)}{library/tree/tree\_dfs.cpp}
needs n, adj, root. Subtree of u = Euler positions [tin[u], tout[u]).
\begin{lstlisting}
// after it: par, depth, sub (subtree size), tin/tout (Euler tour)
// subtree of u = positions [tin[u], tout[u]) -> a range query on the tour
\end{lstlisting}
\end{libentry}
\begin{libentry}{Tree diameter}{Ctrl+Alt+T, Shift+D}{cpdiameter}{above solve()}{O(n)}{library/tree/tree\_diameter.cpp}
auto [len, path] = tree\_diameter(adj); two BFS passes, path lists the vertices end to end.
\begin{lstlisting}
auto [len, path] = tree_diameter(adj);
\end{lstlisting}
\end{libentry}
\begin{libentry}{Tree DP (subtree DP, iterative)}{Ctrl+Alt+T, ;}{cptreedp}{inside solve()}{O(n)}{library/tree/tree\_dp.cpp}
needs n, adj, root. Children before parents (reverse BFS order); example: maximum independent set.
\begin{lstlisting}
// answer: max(dp[root][0], dp[root][1]); change the two lines marked "transition"
\end{lstlisting}
\end{libentry}
\begin{libentry}{Tree isomorphism (canonical IDs, AHU)}{Ctrl+Alt+U, H}{cptreehash}{above solve()}{O(n log n)}{library/tree/tree\_hash.cpp}
TreeCanon tc; tc.rooted(adj, root) and tc.unrooted(adj) give integer IDs: two trees get the same ID (from the same TreeCanon object) iff they are isomorphic. rooted\_all(adj, root)[v] is the ID of v\textquotesingle{}s subtree.
\begin{lstlisting}
TreeCanon tc;
bool iso = tc.unrooted(adj1) == tc.unrooted(adj2);
\end{lstlisting}
\end{libentry}
\begin{libentry}{Tree knapsack (merging child DPs)}{Ctrl+Alt+Y, Shift+K}{cptreeknap}{inside solve()}{O(n\textasciicircum{}2)}{library/tree/tree\_knapsack.cpp}
needs n, adj, root, w. best[v][j] = max value choosing j vertices of v\textquotesingle{}s subtree, each chosen vertex\textquotesingle{}s parent chosen too.
\begin{lstlisting}
// answer for exactly k vertices containing the root: best[root][k]; merging costs O(n^2) in total
\end{lstlisting}
\end{libentry}
\begin{libentry}{Virtual (auxiliary) tree}{Ctrl+Alt+Y, V}{cpvirtual}{above solve()}{O(k log k) LCA calls}{library/tree/virtual\_tree.cpp}
virtual\_tree(nodes, tin, lca) -\textgreater{} \{vertices, parent\}: key nodes + pairwise LCAs, parent[i] = virtual parent of vertices[i] (-1 root).
\begin{lstlisting}
auto [vs, vpar] = virtual_tree(keys, tin, [&](int a, int b) { return t.lca(a, b); });  // tin from cptree, t from cplca
// then DP over vs in reverse order: children appear after their parents
\end{lstlisting}
\end{libentry}
\section{Data structures \& range queries}
\begin{libentry}{Anti-hack hash for unordered\_map}{Ctrl+Alt+T, Shift+H}{cphash}{above solve()}{O(1) expected per operation}{library/ds/custom\_hash.cpp}
unordered\_map\textless{}ll, int, custom\_hash\textgreater{} mp; also works for pair keys. Stops anti-hash TLE hacks.
\begin{lstlisting}
unordered_map<ll, int, custom_hash> cnt;
unordered_map<pair<int, int>, int, custom_hash> seen;
\end{lstlisting}
\end{libentry}
\begin{libentry}{Cartesian tree}{Ctrl+Alt+U, C}{cpcartesian}{above solve()}{O(n) (one stack pass)}{library/ds/cartesian\_tree.cpp}
CartesianTree ct(a): min-heap on values, in-order = array order (ties: the leftmost is higher). ct.root, ct.par[i], ct.lc[i], ct.rc[i] (-1 = none). The subtree of i is the maximal range where a[i] is the minimum.
\begin{lstlisting}
CartesianTree ct(a);
int r = ct.root; for (int i = 0; i < n; i++) { int p = ct.par[i]; }
\end{lstlisting}
\end{libentry}
\begin{libentry}{CDQ divide and conquer (3D dominance counting)}{Ctrl+Alt+U, 3}{cpcdq}{above solve()}{O(n log\textasciicircum{}2 n)}{library/ds/cdq.cpp}
cdq\_dominance(pts): for each point i, the number of j != i with a\_j \textless{}= a\_i, b\_j \textless{}= b\_i, c\_j \textless{}= c\_i. The pattern (sort by one key, divide and conquer on the second, BIT on the third) solves offline 3D problems.
\begin{lstlisting}
vector<array<ll, 3>> pts(n);   // (a, b, c)
vector<int> dom = cdq_dominance(pts);
\end{lstlisting}
\end{libentry}
\begin{libentry}{Disjoint sparse table (O(1) queries for any associative operation)}{Ctrl+Alt+U, T}{cpdisjointsparse}{above solve()}{O(n log n) build, O(1) query}{library/ds/disjoint\_sparse\_table.cpp}
DisjointSparse\textless{}T, Op\textgreater{} t(a, op); t.query(l, r) = a[l] op ... op a[r-1] for [l, r), l \textless{} r. Works for sums, products mod p, matrix products: no need for idempotence.
\begin{lstlisting}
DisjointSparse t(a, [](ll x, ll y) { return x * y % MOD; });
ll prod = t.query(l, r);
\end{lstlisting}
\end{libentry}
\begin{libentry}{Distinct values in ranges (offline sweep + BIT)}{Ctrl+Alt+Y, Shift+D}{cpdistinct}{above solve()}{O((n + q) log n)}{library/ds/distinct\_in\_range.cpp}
distinct\_in\_ranges(a, queries): number of distinct values in a[l, r) per query. Sweeps r, keeps only last occurrences.
\begin{lstlisting}
vector<int> res = distinct_in_ranges(a, qs);  // qs[i] = {l, r}
\end{lstlisting}
\end{libentry}
\begin{libentry}{Dynamic (sparse) segment tree}{Ctrl+Alt+Y, D}{cpdynseg}{above solve()}{O(log range) per operation, O(q log range) memory}{library/ds/dynamic\_segtree.cpp}
DynSeg st(lo, hi): positions in [lo, hi) as large as 1e18, nodes made on demand; st.add(pos, v), st.sum(l, r) on [l, r).
\begin{lstlisting}
DynSeg st(0, 1e9);
st.add(x, 1); ll c = st.sum(l, r);
\end{lstlisting}
\end{libentry}
\begin{libentry}{Euler tour tree (dynamic forest: link, cut, connectivity, component sums)}{Ctrl+Alt+U, E}{cpett}{above solve()}{O(log n) expected per operation}{library/ds/euler\_tour\_tree.cpp}
ETT f(n); f.link(u, v), f.cut(u, v) (the edge must exist), f.connected(u, v), f.set(v, x), f.comp\_sum(v), f.comp\_size(v). Each tree is stored as its Euler tour in a treap.
\begin{lstlisting}
ETT f(n); f.link(0, 1); f.set(1, 5);
bool c = f.connected(0, 1); ll s = f.comp_sum(0); f.cut(0, 1);
\end{lstlisting}
\end{libentry}
\begin{libentry}{Fenwick tree / BIT}{Ctrl+Alt+T, I}{cpbit}{above solve()}{O(log n) per operation}{library/ds/fenwick.cpp}
BIT b(n); b.add(i, v), b.sum(l, r) on [l, r), b.prefix(i) = sum of [0, i), b.lower\_bound(k).
\begin{lstlisting}
BIT b(n);
b.add(i, v); b.sum(l, r); b.prefix(i); b.lower_bound(k);
\end{lstlisting}
\end{libentry}
\begin{libentry}{Fenwick tree 2D}{Ctrl+Alt+Y, Shift+I}{cpbit2d}{above solve()}{O(log n log m) per operation}{library/ds/fenwick\_2d.cpp}
BIT2D b(n, m); b.add(r, c, v); b.sum(r1, c1, r2, c2) = sum over rows [r1, r2) x cols [c1, c2).
\begin{lstlisting}
BIT2D b(n, m);
b.add(r, c, 1); ll s = b.sum(r1, c1, r2, c2);
\end{lstlisting}
\end{libentry}
\begin{libentry}{Fenwick tree: range add + range sum}{Ctrl+Alt+Y, I}{cpbitrange}{above solve()}{O(log n) per operation}{library/ds/fenwick\_range.cpp}
BITRange b(n); b.add(l, r, v) adds v on [l, r); b.sum(l, r). Two Fenwick trees over the difference array.
\begin{lstlisting}
BITRange b(n);
b.add(l, r, v); ll s = b.sum(l, r);
\end{lstlisting}
\end{libentry}
\begin{libentry}{Implicit treap (sequence with split / merge)}{Ctrl+Alt+Y, T}{cptreap}{above solve()}{O(log n) expected per operation}{library/ds/treap.cpp}
Treap tr; tr.insert(pos, v), tr.erase(pos), tr.reverse\_range(l, r), tr.range\_sum(l, r), tr.to\_vector(). Expected O(log n).
\begin{lstlisting}
Treap tr; for (auto x : a) tr.insert(tr.count(), x);
tr.reverse_range(l, r); ll s = tr.range_sum(l, r); vector<ll> now = tr.to_vector();
\end{lstlisting}
\end{libentry}
\begin{libentry}{Lazy segment tree (range add; range sum + min)}{Ctrl+Alt+T, Shift+S}{cplazy}{above solve()}{O(log n) per operation}{library/ds/lazy\_segment\_tree.cpp}
LazySeg st(n) or st(a); st.add(l, r, v), st.sum\_of(l, r), st.min\_of(l, r), all on [l, r).
\begin{lstlisting}
LazySeg st(a);
st.add(l, r, v); st.sum_of(l, r); st.min_of(l, r);  // all [l, r)
\end{lstlisting}
\end{libentry}
\begin{libentry}{Lazy segment tree (range assign + range add; sum, min)}{Ctrl+Alt+Y, A}{cplazyassign}{above solve()}{O(log n) per operation}{library/ds/lazy\_assign.cpp}
LazyAssign st(n or a): st.assign(l, r, v), st.add(l, r, v), st.sum\_of(l, r), st.min\_of(l, r), all on [l, r).
\begin{lstlisting}
LazyAssign st(a);
st.assign(l, r, 5); st.add(l, r, -2); st.sum_of(l, r); st.min_of(l, r);
\end{lstlisting}
\end{libentry}
\begin{libentry}{Mo\textquotesingle{}s algorithm (offline range queries)}{Ctrl+Alt+Y, O}{cpmo}{above solve()}{O((n + q) sqrt n) add/remove calls}{library/ds/mo\_algorithm.cpp}
mo(n, queries, add, remove, answer): queries [l, r); add(i) / remove(i) update the window, answer(qi) records query qi.
\begin{lstlisting}
vector<int> cnt(maxv), res(q); int distinct = 0;
mo(n, qs, [&](int i) { distinct += cnt[a[i]]++ == 0; }, [&](int i) { distinct -= --cnt[a[i]] == 0; }, [&](int qi) { res[qi] = distinct; });
\end{lstlisting}
\end{libentry}
\begin{libentry}{Mo\textquotesingle{}s algorithm on tree paths}{Ctrl+Alt+U, O}{cpmotree}{above solve()}{O((n + q) sqrt(n) * cost of add / remove + q log n)}{library/ds/mo\_tree.cpp}
mo\_on\_tree(adj, root, qs, add, remove, answer): offline queries on the vertex set of the path u-v. add(v)/remove(v) change the current set, answer(i) records query i. Uses the Euler tour trick (a vertex seen twice cancels).
\begin{lstlisting}
mo_on_tree(adj, 0, qs, [&](int v) { if (cnt[c[v]]++ == 0) distinct++; }, [&](int v) { if (--cnt[c[v]] == 0) distinct--; }, [&](int i) { ans[i] = distinct; });
\end{lstlisting}
\end{libentry}
\begin{libentry}{Mo\textquotesingle{}s algorithm with updates (time as a third dimension)}{Ctrl+Alt+U, Shift+O}{cpmoupdates}{above solve()}{O(n\textasciicircum{}(5/3)) for n queries and updates}{library/ds/mo\_updates.cpp}
mo\_with\_updates(n, qs, add, remove, apply, answer): qs[i] = \{l, r, t\} asks about [l, r) after the first t updates. apply(u, L, R) toggles update u (swap in/out) given the window [L, R).
\begin{lstlisting}
mo_with_updates(n, qs, add, remove, [&](int u, int L, int R) { auto& [p, v] = upd[u]; if (L <= p && p < R) remove(p); swap(a[p], v); if (L <= p && p < R) add(p); }, [&](int i) { ans[i] = cur; });
\end{lstlisting}
\end{libentry}
\begin{libentry}{Ordered set (PBDS order statistics)}{Ctrl+Alt+T, Shift+O}{cposet}{above solve()}{O(log n) per operation}{library/ds/ordered\_set.cpp}
ordered\_set\textless{}int\textgreater{} s; s.find\_by\_order(k) = k-th smallest (0-based), s.order\_of\_key(x) = \# elements \textless{} x.
\begin{lstlisting}
ordered_set<int> s;
s.insert(x); *s.find_by_order(k); s.order_of_key(x);
\end{lstlisting}
\end{libentry}
\begin{libentry}{Persistent implicit treap (rope with versions)}{Ctrl+Alt+U, R}{cpptreap}{above solve()}{O(log n) expected per operation, O(log n) new nodes per update}{library/ds/persistent\_treap.cpp}
PTreap t; int v = t.build(a); every operation returns a NEW version and keeps the old: insert(v, pos, x), erase(v, pos), concat(v1, v2), slice(v, l, r). Queries: get(v, pos), sum(v, l, r), size(v).
\begin{lstlisting}
PTreap t; int v0 = t.build(a); int v1 = t.insert(v0, 2, 7);
int copy_paste = t.concat(v1, t.slice(v1, l, r)); ll s = t.sum(v0, l, r);
t.t.reserve(n + 50 * q);  // ~46 new 32-byte nodes per update; regrowing briefly holds two copies
\end{lstlisting}
\end{libentry}
\begin{libentry}{Persistent segment tree (k-th smallest in range)}{Ctrl+Alt+Y, P}{cppersist}{above solve()}{O(log n) per operation and O(log n) new nodes per update}{library/ds/persistent\_segtree.cpp}
PersistentSeg ps(n): root = ps.build(a) in O(n) nodes; root2 = ps.update(root1, pos, v) keeps every version; ps.query(root, l, r); ps.kth(rootL, rootR, k).
\begin{lstlisting}
vector<int> root{0}; for (int i = 0; i < n; i++) root.push_back(ps.update(root[i], rank[i], 1));  // ranks from cpcompress
int idx = ps.kth(root[l], root[r], k);  // k-th smallest (0-based) of a[l, r) is vals[idx]
int r0 = ps.build(a);  // a version holding the values a, instead of n updates (n log n nodes)
\end{lstlisting}
\end{libentry}
\begin{libentry}{Segment tree (point update, range query)}{Ctrl+Alt+T, S}{cpseg}{above solve()}{O(log n) per operation}{library/ds/segment\_tree.cpp}
SegTree st(n) or st(a); st.set(i, x), st.get(i), st.query(l, r) on [l, r). Edit T / ID / op.
\begin{lstlisting}
SegTree st(a);  // or SegTree st(n)
st.set(i, x); st.get(i); st.query(l, r);  // [l, r)
\end{lstlisting}
\end{libentry}
\begin{libentry}{Segment tree beats (range chmin, sum, max)}{Ctrl+Alt+Y, S}{cpbeats}{above solve()}{O((n + q) log\textasciicircum{}2 n) amortized}{library/ds/segtree\_beats.cpp}
Beats st(a): st.chmin(l, r, x) sets a[i] = min(a[i], x) on [l, r); st.sum(l, r), st.max(l, r). Amortized O(log\textasciicircum{}2 n).
\begin{lstlisting}
Beats st(a);
st.chmin(l, r, x); ll s = st.sum(l, r); ll m = st.max(l, r);
\end{lstlisting}
\end{libentry}
\begin{libentry}{Segment tree with custom nodes + binary search}{Ctrl+Alt+Y, N}{cpsegnode}{above solve()}{O(log n) per operation}{library/ds/segtree\_node.cpp}
SegNode st(n or leaves): any associative Node (example: max subarray sum); st.max\_right(l, pred) walks down in O(log n).
\begin{lstlisting}
SegNode st(n); st.set(i, Node::leaf(x)); Node r = st.query(l, r); r.best;
int r = st.max_right(l, [&](const Node& x) { return x.sum < k; });  // first r with sum of [l, r] >= k is r
\end{lstlisting}
\end{libentry}
\begin{libentry}{Segment tree: range add, min + count of min}{Ctrl+Alt+Y, Shift+C}{cpmincount}{above solve()}{O(log n) per operation}{library/ds/min\_count\_segtree.cpp}
MinCount st(weights): add(l, r, v); query(l, r) = \{min, total weight at the min\}; zero\_weight() = weight of cells equal to 0 (values must stay \textgreater{}= 0, as coverage counts do).
\begin{lstlisting}
MinCount st(len);  // len[i] = width of elementary interval i (all 1 for plain counting)
st.add(l, r, +1); st.add(l, r, -1); ll covered = total_width - st.zero_weight();  // union of rectangles sweep
\end{lstlisting}
\end{libentry}
\begin{libentry}{Sparse table (O(1) range min)}{Ctrl+Alt+T, R}{cpsparse}{above solve()}{O(n log n) build, O(1) per query}{library/ds/sparse\_table.cpp}
SparseTable st(a); st.query(l, r) = min of [l, r), l \textless{} r. Swap op for max/gcd (idempotent only).
\begin{lstlisting}
SparseTable st(a);
st.query(l, r);  // min of [l, r)
\end{lstlisting}
\end{libentry}
\begin{libentry}{Wavelet tree (k-th smallest / count \textless{}= x in range)}{Ctrl+Alt+Y, W}{cpwavelet}{above solve()}{O(n log n) build, O(log n) per query}{library/ds/wavelet\_tree.cpp}
Wavelet wt(a): wt.kth(l, r, k) = k-th smallest (0-based) of a[l, r); wt.count\_leq(l, r, x). Any ll values (compressed inside); a bit-packed wavelet matrix, so all levels stay in cache.
\begin{lstlisting}
Wavelet wt(a);
ll med = wt.kth(l, r, (r - l) / 2); int c = wt.count_leq(l, r, x);
\end{lstlisting}
\end{libentry}
\section{Dynamic programming}
\begin{libentry}{(min, +) convolution with convex sequences}{Ctrl+Alt+U, Shift+C}{cpminplus}{above solve()}{O(n + m) / O((n + m) log(n + m))}{library/dp/minplus\_convex.cpp}
c[k] = min over i + j = k of a[i] + b[j]. minplus\_convex\_convex(a, b): both convex, O(n + m) by merging slopes; minplus\_convex\_any(a, b): a convex, b arbitrary, O((n + m) log).
\begin{lstlisting}
vector<ll> c = minplus_convex_convex(a, b);
vector<ll> d = minplus_convex_any(convex_a, arbitrary_b);
\end{lstlisting}
\end{libentry}
\begin{libentry}{Bitmask DP (travelling salesman)}{Ctrl+Alt+T, ]}{cptsp}{above solve()}{O(2\textasciicircum{}n n\textasciicircum{}2)}{library/dp/bitmask\_tsp.cpp}
tsp(dist): shortest cycle through all n \textless{}= \textasciitilde{}18 vertices from 0. dp[mask][v] = best path covering mask, ending at v.
\begin{lstlisting}
ll best = tsp(dist);  // dist[i][j] = cost i -> j
// Hamiltonian path count: same loops with dp[mask | 1 << v][v] += dp[mask][u] over edges
\end{lstlisting}
\end{libentry}
\begin{libentry}{Broken profile DP (domino tilings)}{Ctrl+Alt+Y, Shift+P}{cpprofile}{above solve()}{O(n m 2\textasciicircum{}m)}{library/dp/broken\_profile.cpp}
domino\_tilings(n, m): ways to tile an n x m board with 1 x 2 dominoes mod MOD, cell by cell with an m-bit profile (m \textless{}= \textasciitilde{}12).
\begin{lstlisting}
ll ways = domino_tilings(n, m);  // put the smaller side as m
\end{lstlisting}
\end{libentry}
\begin{libentry}{Digit DP}{Ctrl+Alt+T, =}{cpdigit}{above solve()}{O(digits * states * 10)}{library/dp/digit\_dp.cpp}
count\_upto(N) = how many x in [0, N] have no two equal adjacent digits. Change the state / check for your property.
\begin{lstlisting}
ll ans = count_upto(B) - count_upto(A - 1);  // count in [A, B]
\end{lstlisting}
\end{libentry}
\begin{libentry}{Divide \& conquer DP optimization}{Ctrl+Alt+Y, Q}{cpdnc}{above solve()}{O(K n log n) cost calls}{library/dp/dnc\_optimization.cpp}
dnc\_opt(K, n, cost): min total cost splitting a[0, n) into K groups; cost(j, i) = cost of group a[j, i). Needs monotone optimal split points.
\begin{lstlisting}
ll best = dnc_opt(K, n, [&](int j, int i) { ll s = pre[i] - pre[j]; return s * s; });
\end{lstlisting}
\end{libentry}
\begin{libentry}{Grid paths DP}{Ctrl+Alt+T, Shift+Z}{cpgridpaths}{inside solve()}{O(nm)}{library/dp/grid\_paths.cpp}
needs n, m, grid. ways[r][c] = right/down paths from (0, 0) to (r, c) avoiding \textquotesingle{}\#\textquotesingle{}, mod MOD.
\begin{lstlisting}
// answer: ways[n - 1][m - 1]; for a max-sum path use max(...) + value instead of the sum
\end{lstlisting}
\end{libentry}
\begin{libentry}{Knapsack (0/1, unbounded, coin-change ways)}{Ctrl+Alt+T, Shift+Q}{cpknap}{above solve()}{O(n W)}{library/dp/knapsack.cpp}
knapsack\_01(items, W) / knapsack\_unbounded(items, W): max value, weight \textless{}= W; coin\_ways(coins, X): combinations mod MOD.
\begin{lstlisting}
vector<pair<int, ll>> items;  // {weight, value}
ll best = knapsack_01(items, W); ll ways = coin_ways(coins, X);
\end{lstlisting}
\end{libentry}
\begin{libentry}{Knuth optimization (interval DP)}{Ctrl+Alt+Y, Shift+Q}{cpknuth}{above solve()}{O(n\textasciicircum{}2)}{library/dp/knuth\_optimization.cpp}
knuth\_opt(n, cost): dp[l][r] = min over l \textless{} m \textless{} r of dp[l][m] + dp[m][r] + cost(l, r), dp[i][i+1] = 0. Needs monotone opt.
\begin{lstlisting}
ll best = knuth_opt(n, [&](int l, int r) { return pre[r] - pre[l]; });  // optimal merging of a[0, n)
\end{lstlisting}
\end{libentry}
\begin{libentry}{LCS + edit distance}{Ctrl+Alt+T, [}{cplcs}{above solve()}{O(nm)}{library/dp/lcs\_edit.cpp}
lcs(a, b) -\textgreater{} \{length, one LCS\}; edit\_distance(a, b) = min insert / delete / replace operations.
\begin{lstlisting}
auto [len, s] = lcs(a, b);
int d = edit_distance(a, b);
\end{lstlisting}
\end{libentry}
\begin{libentry}{Li Chao tree (min of lines)}{Ctrl+Alt+Y, Shift+L}{cplichao}{above solve()}{O(log range) per add / query}{library/dp/li\_chao.cpp}
LiChao lc(lo, hi); lc.add(k, m); lc.query(x) = min of k*x + m for integer x in [lo, hi). Nodes created on demand.
\begin{lstlisting}
LiChao lc(-1e9, 1e9 + 1);
lc.add(k, m); ll best = lc.query(x);
\end{lstlisting}
\end{libentry}
\begin{libentry}{LineContainer (convex hull trick, any order)}{Ctrl+Alt+Y, L}{cplinecontainer}{above solve()}{O(log n) per add / query}{library/dp/line\_container.cpp}
LineContainer lc; lc.add(k, m); lc.query(x) = max of k*x + m. For a minimum add(-k, -m) and negate the query.
\begin{lstlisting}
LineContainer lc;
lc.add(slope, intercept); ll best = lc.query(x);  // dp[i] = max_j (dp[j] + b[j] * a[i]) in O(log n)
\end{lstlisting}
\end{libentry}
\begin{libentry}{Longest increasing subsequence}{Ctrl+Alt+T, Shift+L}{cplis}{above solve()}{O(n log n)}{library/dp/lis.cpp}
vector\textless{}int\textgreater{} idx = lis(a): indices of one strictly increasing LIS, O(n log n). Length = idx.size().
\begin{lstlisting}
vector<int> idx = lis(a);  // idx.size() = LIS length
\end{lstlisting}
\end{libentry}
\begin{libentry}{Monotone convex hull trick (sorted slopes)}{Ctrl+Alt+Y, Z}{cpmonocht}{above solve()}{O(1) amortized add, O(log n) query, O(1) amortized query\_inc}{library/dp/monotone\_cht.cpp}
MonoCHT h; h.add(k, m) with non-increasing k; h.query(x) = min of k*x + m (binary search); h.query\_inc(x) for non-decreasing x (pointer walk). For max: add(-k, -m), negate.
\begin{lstlisting}
MonoCHT h; h.add(b[j], dp[j]);
dp[i] = h.query_inc(a[i]) + c[i];
\end{lstlisting}
\end{libentry}
\begin{libentry}{Monotone minima (row minima when the argmin moves right)}{Ctrl+Alt+U, M}{cpmonominima}{above solve()}{O((n + m) log n) evaluations of f}{library/dp/monotone\_minima.cpp}
monotone\_minima(n, m, f): for each row i, the leftmost column j minimizing f(i, j), given that this argmin is non-decreasing in i (e.g. a Monge matrix).
\begin{lstlisting}
vector<int> opt = monotone_minima(n, m, [&](int i, int j) { return cost(i, j); });
\end{lstlisting}
\end{libentry}
\begin{libentry}{Range (interval) DP}{Ctrl+Alt+T, -}{cprangedp}{inside solve()}{O(n\textasciicircum{}3)}{library/dp/range\_dp.cpp}
needs n, a. dp[l][r] over intervals by increasing length; example: merge adjacent piles, cost = merged sum.
\begin{lstlisting}
// answer: dp[0][n - 1]; change the inner transition for your problem
\end{lstlisting}
\end{libentry}
\begin{libentry}{Slope trick (make an array non-decreasing)}{Ctrl+Alt+Y, Shift+S}{cpslope}{above solve()}{O(n log n)}{library/dp/slope\_trick.cpp}
make\_non\_decreasing(a): min total \textbar{}a[i] - b[i]\textbar{} over non-decreasing b. Max-heap of slope breakpoints. Strictly increasing: use a[i] - i.
\begin{lstlisting}
ll cost = make_non_decreasing(a);
\end{lstlisting}
\end{libentry}
\begin{libentry}{Subset convolution}{Ctrl+Alt+U, Shift+S}{cpsubsetconv}{above solve()}{O(n\textasciicircum{}2 2\textasciicircum{}n)}{library/dp/subset\_convolution.cpp}
subset\_convolution(a, b): c[S] = sum over T subset of S of a[T] * b[S \textbackslash{} T] (disjoint unions only), mod MOD. Sizes 2\textasciicircum{}n; n = 20 takes about 1 s and 176 MB.
\begin{lstlisting}
vector<ll> c = subset_convolution(a, b);   // a.size() == b.size() == 1 << n
\end{lstlisting}
\end{libentry}
\begin{libentry}{Subset sums with bitset}{Ctrl+Alt+T, Shift+1}{cpbitsetsum}{above solve()}{O(n MAXS / 64)}{library/dp/subset\_sum\_bitset.cpp}
subset\_sums\textless{}MAXS\textgreater{}(a): bitset with bit s set iff some subset of a sums to s (s \textless{} MAXS). 64x faster than a bool DP.
\begin{lstlisting}
auto can = subset_sums<100001>(a);
if (can[target]) { ... }
\end{lstlisting}
\end{libentry}
\begin{libentry}{Sum over subsets (SOS DP)}{Ctrl+Alt+Y, U}{cpsos}{above solve()}{O(B 2\textasciicircum{}B)}{library/dp/sos\_dp.cpp}
sos\_subsets(f, B): f[mask] becomes the sum of f[sub] over all submasks; sos\_supersets: over supermasks. In place, size 2\textasciicircum{}B.
\begin{lstlisting}
vector<ll> f(1 << B); /* f[mask] = count of numbers equal to mask */ sos_subsets(f, B);
// f[mask] now counts numbers x with (x & mask) == x
\end{lstlisting}
\end{libentry}
\section{Math}
\begin{libentry}{Berlekamp-Massey (find a linear recurrence)}{Ctrl+Alt+U, Shift+B}{cpbm}{above solve()}{O(n\textasciicircum{}2)}{library/math/berlekamp\_massey.cpp}
berlekamp\_massey(s): shortest c with s[i] = c[0]*s[i-1] + c[1]*s[i-2] + ... mod MOD (prime). Feed 2L terms to find a recurrence of order L.
\begin{lstlisting}
vector<ll> c = berlekamp_massey({1, 1, 2, 3, 5, 8, 13, 21});   // {1, 1}
ll x = linear_rec_nth(first_terms, c, n);   // with the Kitamasa template
\end{lstlisting}
\end{libentry}
\begin{libentry}{Chinese remainder theorem (any moduli)}{Ctrl+Alt+Y, 1}{cpcrt}{above solve()}{O(log m)}{library/math/crt.cpp}
crt(r1, m1, r2, m2) -\textgreater{} \{r, lcm\}: x = r (mod lcm) solves both congruences; \{-1, -1\} if none. Moduli need not be coprime.
\begin{lstlisting}
auto [r, m] = crt(2, 3, 3, 5);  // x = 8 (mod 15)
// many congruences: fold left, feeding each result back in as (r1, m1); lcm must stay < ~4e18
\end{lstlisting}
\end{libentry}
\begin{libentry}{Determinant mod p, matrix-tree theorem}{Ctrl+Alt+U, Shift+D}{cpdet}{above solve()}{O(n\textasciicircum{}3)}{library/math/determinant.cpp}
det\_mod(A): determinant mod MOD (prime) by Gaussian elimination; count\_spanning\_trees(n, edges): Kirchhoff\textquotesingle{}s matrix-tree theorem (multi-edges allowed). LGV lemma: det of path-count matrix.
\begin{lstlisting}
ll d = det_mod(A);
ll t = count_spanning_trees(n, edges);   // edges: 0-indexed pairs
\end{lstlisting}
\end{libentry}
\begin{libentry}{Discrete logarithm (baby-step giant-step)}{Ctrl+Alt+U, L}{cpdlog}{above solve()}{O(sqrt(m)) time and memory}{library/math/discrete\_log.cpp}
discrete\_log(a, b, m): smallest x \textgreater{}= 0 with a\textasciicircum{}x = b (mod m), or -1. Any m (not only primes), m \textless{}= 3e9.
\begin{lstlisting}
ll x = discrete_log(2, 5, 13);   // 2^9 = 512 = 5 (mod 13) -> 9
\end{lstlisting}
\end{libentry}
\begin{libentry}{Divisor sieve (count + sum of divisors)}{Ctrl+Alt+T, Shift+2}{cpdivsieve}{above solve()}{O(N log N)}{library/math/divisor\_sieve.cpp}
divisor\_sieve(N) -\textgreater{} \{cnt, sum\}: number and sum of divisors of every x \textless{}= N by looping over multiples.
\begin{lstlisting}
auto [cnt, sum] = divisor_sieve(1000000);
\end{lstlisting}
\end{libentry}
\begin{libentry}{Du\textquotesingle{}s sieve: prefix sums of phi and mu up to 1e10}{Ctrl+Alt+U, D}{cpdusieve}{above solve()}{O(n\textasciicircum{}(2/3)) for the first query, memoised after}{library/math/du\_sieve.cpp}
DuSieve ds(n\_max); ds.mu\_sum(n) = mu(1)+...+mu(n) (exact), ds.phi\_sum(n) = phi(1)+...+phi(n) mod MOD.
\begin{lstlisting}
DuSieve ds(10'000'000'000LL);
ll M = ds.mu_sum(n), Phi = ds.phi_sum(n);
\end{lstlisting}
\end{libentry}
\begin{libentry}{Euler\textquotesingle{}s totient}{Ctrl+Alt+Y, Shift+E}{cpphi}{above solve()}{O(sqrt n) single, O(N log log N) sieve}{library/math/totient.cpp}
phi(n) by trial division; phi\_sieve(N) = phi of every x \textless{}= N. phi(n) counts 1 \textless{}= k \textless{}= n with gcd(k, n) = 1.
\begin{lstlisting}
ll p = phi(n);
auto ph = phi_sieve(1000000);  // a^phi(m) = 1 (mod m) when gcd(a, m) = 1
\end{lstlisting}
\end{libentry}
\begin{libentry}{Extended GCD + inverse for any modulus}{Ctrl+Alt+T, X}{cpextgcd}{above solve()}{O(log min(a, b))}{library/math/ext\_gcd.cpp}
g = ext\_gcd(a, b, x, y) with a*x + b*y = g; inv\_mod(a, m) for gcd(a, m) = 1, else -1.
\begin{lstlisting}
ll x, y, g = ext_gcd(a, b, x, y);  // a*x + b*y = g
inv_mod(a, m);  // -1 if gcd(a, m) != 1
\end{lstlisting}
\end{libentry}
\begin{libentry}{Factorize + divisors (trial division)}{Ctrl+Alt+T, Shift+P}{cpfactor}{above solve()}{O(sqrt x)}{library/math/factorize.cpp}
factorize(x) -\textgreater{} \{prime, exponent\}, divisors(x) sorted. O(sqrt x), fine up to \textasciitilde{}1e12.
\begin{lstlisting}
for (auto [p, e] : factorize(x)) { ... }
vector<ll> d = divisors(x);
\end{lstlisting}
\end{libentry}
\begin{libentry}{FFT (exact integer convolution)}{Ctrl+Alt+Y, F}{cpfft}{above solve()}{O((n + m) log(n + m))}{library/math/fft.cpp}
multiply(a, b): convolution c[k] = sum a[i] b[k - i] with complex FFT; exact while results stay below \textasciitilde{}1e15.
\begin{lstlisting}
vector<ll> c = multiply(a, b);  // polynomial product, counting pairs with a given sum, ...
\end{lstlisting}
\end{libentry}
\begin{libentry}{Floor-division blocks}{Ctrl+Alt+Y, Shift+F}{cpfloorblocks}{above solve()}{O(sqrt n)}{library/math/floor\_blocks.cpp}
for\_each\_floor\_block(n, f): f(l, r, q) for every maximal [l, r] where n / i == q; only O(sqrt n) blocks.
\begin{lstlisting}
ll s = 0; for_each_floor_block(n, [&](ll l, ll r, ll q) { s += q * (r - l + 1); });  // sum of n / i
// sum of divisors of 1..n: q * (l + r) * (r - l + 1) / 2 per block
\end{lstlisting}
\end{libentry}
\begin{libentry}{Gaussian elimination mod p (linear systems, rank)}{Ctrl+Alt+U, N}{cpgauss}{above solve()}{O(n m min(n, m))}{library/math/gauss\_mod.cpp}
gauss\_mod(a, b, x): solves a x = b (mod MOD, MOD prime) for an n x m matrix a. Returns the rank, or -1 if there is no solution; x receives one solution (free variables set to 0).
\begin{lstlisting}
vector<vector<ll>> a(n, vector<ll>(m)); vector<ll> b(n), x;
int r = gauss_mod(a, b, x);   // -1: inconsistent; else m - r free variables, MOD^(m - r) solutions
\end{lstlisting}
\end{libentry}
\begin{libentry}{Grundy numbers (Sprague-Grundy)}{Ctrl+Alt+Y, G}{cpgrundy}{above solve()}{O(N * moves)}{library/math/grundy.cpp}
grundy\_numbers(N, moves): g[s] = mex of g over moves(s) (moves must go to smaller states). XOR of games != 0 -\textgreater{} first player wins.
\begin{lstlisting}
auto g = grundy_numbers(N, [&](int s) { vector<int> to; for (int k : {1, 3, 4}) if (s >= k) to.push_back(s - k); return to; });
int x = 0; for (int pile : piles) x ^= g[pile]; bool first_wins = x != 0;
\end{lstlisting}
\end{libentry}
\begin{libentry}{Inclusion-exclusion (divisible by any)}{Ctrl+Alt+Y, 3}{cpincexc}{above solve()}{O(2\textasciicircum{}k k)}{library/math/inclusion\_exclusion.cpp}
count\_divisible\_by\_any(N, ds): how many x in [1, N] are divisible by at least one of ds (k \textless{}= \textasciitilde{}20). Odd subsets add, even subtract.
\begin{lstlisting}
ll c = count_divisible_by_any(N, {2, 3, 5});
// coprime to m: N - count_divisible_by_any(N, distinct primes of m)
\end{lstlisting}
\end{libentry}
\begin{libentry}{Lagrange interpolation}{Ctrl+Alt+U, I}{cplagrange}{above solve()}{O(k) consecutive points (plus O(log MOD)), O(k\textasciicircum{}2) general points}{library/math/lagrange.cpp}
lagrange\_consecutive(y, x): the degree \textless{} k polynomial with P(i) = y[i] for i = 0..k-1, evaluated at any x; lagrange\_general(xs, ys, x) for arbitrary distinct points. Mod MOD (prime).
\begin{lstlisting}
ll s = lagrange_consecutive({0, 1, 5, 14}, 1'000'000'000);   // 0^2 + 1^2 + ... + 10^9^2 mod MOD
ll v = lagrange_general(xs, ys, x);
\end{lstlisting}
\end{libentry}
\begin{libentry}{Matrix exponentiation}{Ctrl+Alt+T, Shift+X}{cpmatrix}{above solve()}{O(k\textasciicircum{}3 log e)}{library/math/matrix.cpp}
Mat R = mat\_pow(A, e), mat\_mul(A, B), entries mod MOD. Linear recurrences in O(k\textasciicircum{}3 log e).
\begin{lstlisting}
Mat M = {{1, 1}, {1, 0}};
ll fib_n = mat_pow(M, n)[0][1];
\end{lstlisting}
\end{libentry}
\begin{libentry}{Miller-Rabin + Pollard rho (numbers up to 1e18)}{Ctrl+Alt+Y, X}{cppollard}{above solve()}{\textasciitilde{}O(n\textasciicircum{}(1/4) log n) per factorization}{library/math/pollard\_rho.cpp}
BigFactor::is\_prime(n) (deterministic below 2\textasciicircum{}64), BigFactor::factor(n) -\textgreater{} prime factors with repeats, unsorted.
\begin{lstlisting}
bool p = BigFactor::is_prime(n);
auto f = BigFactor::factor(n); sort(f.begin(), f.end());
\end{lstlisting}
\end{libentry}
\begin{libentry}{Mint (modular integer)}{Ctrl+Alt+T, Shift+M}{cpmint}{above solve()}{O(1) per op, O(log MOD) for / and inv}{library/math/mint.cpp}
mint x = 5; x += 3; x *= x.inv(); x.pow(e); cout \textless{}\textless{} x. Uses the template\textquotesingle{}s MOD (must be prime for /).
\begin{lstlisting}
mint x = 5;
x = x * 3 - 1; x /= 2; x.pow(10); cout << x;
\end{lstlisting}
\end{libentry}
\begin{libentry}{Mobius function sieve}{Ctrl+Alt+Y, Shift+M}{cpmobius}{above solve()}{O(N)}{library/math/mobius.cpp}
mobius\_sieve(N) = mu[0..N] by a linear sieve. Key identity: sum over d \textbar{} g of mu(d) = [g == 1] (count coprime pairs, etc.).
\begin{lstlisting}
auto mu = mobius_sieve(N);
// pairs (a, b) <= N with gcd 1: sum over d of mu[d] * (N / d) * (N / d)
\end{lstlisting}
\end{libentry}
\begin{libentry}{Modular power + inverse}{Ctrl+Alt+T, M}{cpmodpow}{above solve()}{O(log e)}{library/math/modpow.cpp}
modpow(a, e, m = MOD), modinv(a, m = MOD) for prime m. Uses the template\textquotesingle{}s MOD.
\begin{lstlisting}
modpow(2, n);  // 2^n mod MOD
modinv(x);    // x^-1 mod MOD
\end{lstlisting}
\end{libentry}
\begin{libentry}{N-th term of a linear recurrence (Kitamasa)}{Ctrl+Alt+U, K}{cpkitamasa}{above solve()}{O(L\textasciicircum{}2 log n)}{library/math/linear\_recurrence.cpp}
linear\_rec\_nth(a, c, n): a[n] for a[i] = c[0]*a[i-1] + ... + c[L-1]*a[i-L], given a[0..L-1]; mod MOD. Pairs with Berlekamp-Massey.
\begin{lstlisting}
ll f = linear_rec_nth({0, 1}, {1, 1}, 1'000'000'000'000'000'000LL);   // Fibonacci
\end{lstlisting}
\end{libentry}
\begin{libentry}{nCr / factorials mod p}{Ctrl+Alt+T, N}{cpncr}{above solve()}{O(N) init, O(1) per query}{library/math/ncr.cpp}
init\_comb(N) once, then nCr, nPr, catalan, stars\_and\_bars, fact[i], inv\_fact[i]. Uses the template\textquotesingle{}s MOD.
\begin{lstlisting}
init_comb(200000);
nCr(n, r); nPr(n, r); catalan(n); fact[n];
\end{lstlisting}
\end{libentry}
\begin{libentry}{NTT (convolution mod 998244353)}{Ctrl+Alt+Y, Shift+T}{cpntt}{above solve()}{O((n + m) log(n + m))}{library/math/ntt.cpp}
ntt\_multiply(a, b): convolution with every coefficient mod 998244353, exact for any size up to 2\textasciicircum{}23.
\begin{lstlisting}
vector<ll> c = ntt_multiply(a, b);
\end{lstlisting}
\end{libentry}
\begin{libentry}{Pascal triangle (nCr for any modulus)}{Ctrl+Alt+Y, Shift+N}{cppascal}{above solve()}{O(N\textasciicircum{}2)}{library/math/pascal.cpp}
pascal(N, m)[n][r] = C(n, r) mod m for n \textless{}= N; works when m is not prime or is 1e18-ish (then keep N small).
\begin{lstlisting}
auto C = pascal(60, 1e18);  // exact binomials up to C(60, 30) ~ 1.2e17
\end{lstlisting}
\end{libentry}
\begin{libentry}{Polynomial toolkit mod 998244353 (inverse, divmod, log, exp, sqrt, pow, multipoint)}{Ctrl+Alt+U, Shift+P}{cppoly}{above solve()}{O(n log n) for mul / inv / divmod / log / exp / sqrt / pow; O(n log\textasciicircum{}2 n) for multipoint and interpolation}{library/math/poly.cpp}
Poly = vector\textless{}ll\textgreater{}. poly\_mul, poly\_inv(a, n), poly\_divmod(a, b), poly\_log / poly\_exp / poly\_sqrt (a, n), poly\_pow(a, k, n), poly\_eval\_many(a, xs), poly\_interpolate(xs, ys). Self-contained NTT.
\begin{lstlisting}
Poly b = poly_inv(a, n);   // a * b = 1 mod x^n, needs a[0] != 0
auto [q, r] = poly_divmod(a, b);  Poly e = poly_exp(a, n);  vector<ll> v = poly_eval_many(a, xs);
\end{lstlisting}
\end{libentry}
\begin{libentry}{Prime counting and prime sums up to 1e11 (Lucy)}{Ctrl+Alt+U, P}{cpprimecount}{above solve()}{O(n\textasciicircum{}(3/4)) time, O(sqrt n) memory}{library/math/prime\_count.cpp}
prime\_count(n) = number of primes \textless{}= n (exact); prime\_sum(n) = sum of primes \textless{}= n mod MOD.
\begin{lstlisting}
ll pi = prime_count(100'000'000'000LL);   // 4118054813
ll s = prime_sum(n);
\end{lstlisting}
\end{libentry}
\begin{libentry}{Primitive root}{Ctrl+Alt+U, G}{cpprimroot}{above solve()}{O(sqrt(m) + answer * (number of prime factors of phi) * log m)}{library/math/primitive\_root.cpp}
primitive\_root(m): smallest generator g of (Z/mZ)* (exists for m = 1, 2, 4, p\textasciicircum{}k, 2p\textasciicircum{}k), else -1. m \textless{}= 3e9.
\begin{lstlisting}
ll g = primitive_root(998244353);   // 3
\end{lstlisting}
\end{libentry}
\begin{libentry}{Sieve (primes + smallest prime factor)}{Ctrl+Alt+T, P}{cpsieve}{above solve()}{O(N) init, O(log x) per factorization}{library/math/sieve.cpp}
init\_sieve(N); primes, spf[x], is\_prime(x), factor\_small(x) -\textgreater{} \{p, e\} in O(log x) for x \textless{}= N.
\begin{lstlisting}
init_sieve(10000000);
is_prime(x); primes; factor_small(x);  // x <= N
\end{lstlisting}
\end{libentry}
\begin{libentry}{Square root modulo a prime (Tonelli-Shanks)}{Ctrl+Alt+U, Q}{cpsqrtmod}{above solve()}{O(log\textasciicircum{}2 p)}{library/math/sqrt\_mod.cpp}
sqrt\_mod(a, p): some x with x*x = a (mod p) for prime p (the other root is p - x), or -1 if a is not a square.
\begin{lstlisting}
ll x = sqrt_mod(10, 13);   // 6 or 7: 6*6 = 36 = 10 (mod 13)
\end{lstlisting}
\end{libentry}
\begin{libentry}{Stern-Brocot: simplest fraction in an interval, best approximation}{Ctrl+Alt+U, B}{cpsternbrocot}{above solve()}{O(log max) (a Euclid-like walk down the Stern-Brocot tree)}{library/math/stern\_brocot.cpp}
simplest\_between(a, b, c, d): p/q with the smallest q (then p) and a/b \textless{} p/q \textless{} c/d (0 \textless{}= a/b \textless{} c/d); best\_approx(x, y, N): p/q closest to x/y with q \textless{}= N.
\begin{lstlisting}
auto [p, q] = simplest_between(1, 3, 1, 2);   // 2/5
auto [p2, q2] = best_approx(314159, 100000, 100);   // 311/99
\end{lstlisting}
\end{libentry}
\begin{libentry}{Stirling numbers (both kinds) and Bell numbers}{Ctrl+Alt+U, S}{cpstirling}{above solve()}{O(N\textasciicircum{}2) time and memory}{library/math/stirling.cpp}
init\_stirling(N): S2[n][k] (partitions of n items into k nonempty sets), S1[n][k] (permutations of n with k cycles), bell[n]; all mod MOD.
\begin{lstlisting}
init_stirling(1000);
S2[n][k]; S1[n][k]; bell[n];
\end{lstlisting}
\end{libentry}
\begin{libentry}{Walsh-Hadamard transform: XOR / AND / OR convolution}{Ctrl+Alt+U, W}{cpfwht}{above solve()}{O(n log n) for n = padded size}{library/math/fwht.cpp}
xor\_convolution(a, b), and\_convolution(a, b), or\_convolution(a, b): c[k] = sum of a[i]*b[j] over i\textasciicircum{}j = k (i\&j, i\textbar{}j) mod MOD. Sizes padded to a power of two.
\begin{lstlisting}
vector<ll> c = xor_convolution(a, b);
\end{lstlisting}
\end{libentry}
\begin{libentry}{XOR basis (linear basis over GF(2))}{Ctrl+Alt+Y, Shift+X}{cpxorbasis}{above solve()}{O(B) per operation, B = 62 bits}{library/math/xor\_basis.cpp}
XorBasis xb; xb.insert(x) (false if x was already representable); xb.can\_make(x); xb.max\_xor(); 2\textasciicircum{}xb.dim distinct subset XORs.
\begin{lstlisting}
XorBasis xb; for (auto x : a) xb.insert(x);
ll best = xb.max_xor(); bool ok = xb.can_make(target);
\end{lstlisting}
\end{libentry}
\section{Strings}
\begin{libentry}{Aho-Corasick (many patterns at once)}{Ctrl+Alt+Y, Shift+A}{cpaho}{above solve()}{O(total pattern length * 26 + \textbar{}text\textbar{})}{library/string/aho\_corasick.cpp}
AhoCorasick ac; id = ac.add(p) per pattern; ac.build(); ac.occurrences(text)[id] = matches of that pattern. a-z only.
\begin{lstlisting}
AhoCorasick ac; vector<int> id; for (auto& p : pats) id.push_back(ac.add(p)); ac.build();
auto occ = ac.occurrences(text); ll hits = occ[id[i]];  // or walk: s = ac.go(s, ch), ac.out[s] = matches ending here
\end{lstlisting}
\end{libentry}
\begin{libentry}{Binary trie (maximum XOR)}{Ctrl+Alt+T, \textasciigrave{}}{cpxortrie}{above solve()}{O(B) per operation, B = 30 (set B = 61 for 64-bit values)}{library/string/xor\_trie.cpp}
XorTrie t; t.insert(x), t.erase(x) (x must be present), t.max\_xor(x) = max over stored y of x \textasciicircum{} y (non-empty). 30 bits.
\begin{lstlisting}
XorTrie t; t.insert(0);  // prefix XORs: best subarray XOR = max over i of t.max_xor(pref[i])
ll best = t.max_xor(x);
\end{lstlisting}
\end{libentry}
\begin{libentry}{KMP (prefix function + pattern search)}{Ctrl+Alt+T, Shift+K}{cpkmp}{above solve()}{O(n + m)}{library/string/kmp.cpp}
prefix\_function(s) -\textgreater{} pi; find\_occurrences(text, pat) -\textgreater{} start index of every match.
\begin{lstlisting}
vector<int> starts = find_occurrences(text, pat);
vector<int> pi = prefix_function(s);
\end{lstlisting}
\end{libentry}
\begin{libentry}{Lyndon factorization (Duval) and minimal rotation}{Ctrl+Alt+U, Y}{cplyndon}{above solve()}{O(n)}{library/string/lyndon.cpp}
duval(s): start indices of the Lyndon factorization s = w1 w2 ... wk with w1 \textgreater{}= w2 \textgreater{}= ... (each wi strictly smaller than all its rotations); min\_rotation(s): start index of the lexicographically smallest rotation.
\begin{lstlisting}
vector<int> starts = duval(s);
int k = min_rotation(s); string canon = s.substr(k) + s.substr(0, k);
\end{lstlisting}
\end{libentry}
\begin{libentry}{Manacher (all palindromes)}{Ctrl+Alt+Y, 6}{cpmanacher}{above solve()}{O(n)}{library/string/manacher.cpp}
manacher(s) -\textgreater{} \{d1, d2\}: d1[i] = odd palindromes centred at i, d2[i] = even ones centred before i; longest\_palindrome(s) -\textgreater{} \{start, len\}.
\begin{lstlisting}
auto [start, len] = longest_palindrome(s);
// number of palindromic substrings = sum(d1) + sum(d2)
\end{lstlisting}
\end{libentry}
\begin{libentry}{Palindromic tree (eertree)}{Ctrl+Alt+U, Shift+E}{cpeertree}{above solve()}{O(n * alphabet) memory, O(n) amortized time}{library/string/eertree.cpp}
Eertree t; t.add(c) for each char (lowercase). Nodes 2.. are the distinct palindromic substrings (len, link = longest proper palindromic suffix); t.last = longest palindromic suffix of the current prefix. occurrences() counts each palindrome\textquotesingle{}s occurrences.
\begin{lstlisting}
Eertree t; for (char c : s) t.add(c);
int distinct = t.t.size() - 2; vector<ll> occ = t.occurrences();
\end{lstlisting}
\end{libentry}
\begin{libentry}{String hashing (mod 2\textasciicircum{}61-1)}{Ctrl+Alt+T, H}{cpstrhash}{above solve()}{O(n) build, O(1) per substring}{library/string/string\_hash.cpp}
StrHash h(s); h.get(l, r) = hash of s[l, r). Random base, hack-safe; same base for every string.
\begin{lstlisting}
StrHash h(s);
if (h.get(l1, r1) == h.get(l2, r2)) { /* s[l1, r1) == s[l2, r2) */ }
\end{lstlisting}
\end{libentry}
\begin{libentry}{Suffix array + LCP}{Ctrl+Alt+Y, 4}{cpsa}{above solve()}{O(n log n)}{library/string/suffix\_array.cpp}
suffix\_array(s): sa[i] = start of the i-th smallest suffix; lcp\_array(s, sa)[i] = LCP of suffixes sa[i] and sa[i + 1].
\begin{lstlisting}
auto sa = suffix_array(s); auto lcp = lcp_array(s, sa);
// distinct substrings = n(n+1)/2 - sum(lcp)
\end{lstlisting}
\end{libentry}
\begin{libentry}{Suffix automaton}{Ctrl+Alt+Y, 5}{cpsam}{above solve()}{O(n * 26)}{library/string/suffix\_automaton.cpp}
SuffixAutomaton sam(s): smallest automaton of all substrings. sam.count\_distinct(), sam.contains(t); states hold len / link / next.
\begin{lstlisting}
SuffixAutomaton sam(s);
ll d = sam.count_distinct(); bool in = sam.contains(t);
\end{lstlisting}
\end{libentry}
\begin{libentry}{Trie}{Ctrl+Alt+T, Shift+T}{cptrie}{above solve()}{O(\textbar{}s\textbar{}) per operation}{library/string/trie.cpp}
Trie t; t.insert(s); t.count(s) = copies of s; t.prefix\_count(p) = words starting with p. a-z only.
\begin{lstlisting}
Trie t;
t.insert(s); t.count(s); t.prefix_count(p);
\end{lstlisting}
\end{libentry}
\begin{libentry}{Z-function}{Ctrl+Alt+T, Z}{cpz}{above solve()}{O(n)}{library/string/z\_function.cpp}
z\_function(s): z[i] = longest common prefix of s and s[i..]; z[0] = 0.
\begin{lstlisting}
vector<int> z = z_function(pat + '#' + text);  // z[i] == |pat| -> match
\end{lstlisting}
\end{libentry}
\section{Geometry}
\begin{libentry}{Closest pair of points}{Ctrl+Alt+Y, 9}{cpclosest}{above solve()}{O(n log n)}{library/geometry/closest\_pair.cpp}
closest\_pair(pts) = smallest squared distance between two of the points (n \textgreater{}= 2), sweep by x with a set by y. Needs cppoint.
\begin{lstlisting}
ll d2 = closest_pair(pts);
\end{lstlisting}
\end{libentry}
\begin{libentry}{Convex hull (monotone chain)}{Ctrl+Alt+T, Shift+V}{cphull}{above solve()}{O(n log n)}{library/geometry/convex\_hull.cpp}
convex\_hull(pts) -\textgreater{} counter-clockwise hull without collinear points. Pt = pair\textless{}ll, ll\textgreater{}.
\begin{lstlisting}
vector<Pt> hull = convex_hull(pts);
\end{lstlisting}
\end{libentry}
\begin{libentry}{Geometry primitives (integer points)}{Ctrl+Alt+T, Shift+3}{cppoint}{above solve()}{O(1)}{library/geometry/point.cpp}
Point \{x, y\} with + - * dot cross norm2; orient(a, b, c) = +1 left turn, -1 right, 0 collinear; on\_segment(p, a, b). Exact.
\begin{lstlisting}
Point a{1, 2}, b, c; cin >> b >> c;
int o = orient(a, b, c); bool on = on_segment(p, a, b); ll d2 = (a - b).norm2();
\end{lstlisting}
\end{libentry}
\begin{libentry}{Half-plane intersection}{Ctrl+Alt+U, 4}{cphalfplane}{above solve()}{O(n log n)}{library/geometry/half\_plane.cpp}
half\_plane\_intersection(hs, B): hs[i] = \{ax, ay, bx, by\} keeps the points LEFT of the directed line a-\textgreater{}b. Returns the intersection polygon (CCW, long double) clipped to the box [-B, B]\textasciicircum{}2, or empty if it is empty.
\begin{lstlisting}
vector<array<long double, 4>> hs = {{0, 0, 1, 0}, {1, 0, 0, 1}, {0, 1, 0, 0}};   // a triangle
auto poly = half_plane_intersection(hs, 1e9);
\end{lstlisting}
\end{libentry}
\begin{libentry}{Minimum enclosing circle (Welzl, randomized)}{Ctrl+Alt+U, Shift+M}{cpmincircle}{above solve()}{O(n) expected}{library/geometry/min\_enclosing\_circle.cpp}
min\_enclosing\_circle(pts): the smallest circle \{cx, cy, r\} containing all points (long double). Random shuffle + incremental construction.
\begin{lstlisting}
auto [cx, cy, r] = min_enclosing_circle(pts);   // pts: vector<pair<long double, long double>>
\end{lstlisting}
\end{libentry}
\begin{libentry}{Polygon: area, point inside, lattice points}{Ctrl+Alt+Y, 8}{cppolygon}{above solve()}{O(n)}{library/geometry/polygon.cpp}
area2(poly) = 2 * signed area (\textgreater{} 0 counter-clockwise); in\_polygon(poly, p): 1 inside, 0 boundary, -1 outside; boundary\_points. Needs cppoint.
\begin{lstlisting}
ll twice = llabs(area2(poly)); int where = in_polygon(poly, p);
ll interior = (twice - boundary_points(poly) + 2) / 2;  // Pick's theorem
\end{lstlisting}
\end{libentry}
\begin{libentry}{Rectangles (intersection, union of two)}{Ctrl+Alt+T, Shift+4}{cprect}{above solve()}{O(1)}{library/geometry/rectangles.cpp}
Rect \{x1, y1, x2, y2\} (x1 \textless{} x2, y1 \textless{} y2): area(), intersect\_area(a, b), union\_area(a, b). USACO bronze staples.
\begin{lstlisting}
Rect a{0, 0, 4, 3}, b{2, 1, 6, 5};
ll both = intersect_area(a, b), either = union_area(a, b);
\end{lstlisting}
\end{libentry}
\begin{libentry}{Rotating calipers (diameter and width of a convex polygon)}{Ctrl+Alt+U, 5}{cpcalipers}{above solve()}{O(n)}{library/geometry/rotating\_calipers.cpp}
For a convex polygon h (CCW, no three collinear, e.g. from the convex hull template): hull\_diameter2(h) = max squared distance between two vertices (exact); hull\_width(h) = min distance between two parallel lines enclosing it.
\begin{lstlisting}
ll d2 = hull_diameter2(hull);   // hull: vector<pair<ll, ll>>
long double w = hull_width(hull);
\end{lstlisting}
\end{libentry}
\begin{libentry}{Segment intersection}{Ctrl+Alt+Y, 7}{cpseginter}{above solve()}{O(1)}{library/geometry/segment\_intersection.cpp}
segments\_intersect(a, b, c, d): closed segments ab and cd share a point (touching and overlapping count). Needs cppoint.
\begin{lstlisting}
bool hit = segments_intersect(a, b, c, d);
\end{lstlisting}
\end{libentry}
\section{Flows \& matching}
\begin{libentry}{Bipartite matching (Hopcroft-Karp)}{Ctrl+Alt+Y, Shift+H}{cphopcroft}{above solve()}{O(E sqrt V)}{library/flow/hopcroft\_karp.cpp}
HopcroftKarp hk(nl, nr); hk.add\_edge(u, v) with u on the left, v on the right; hk.max\_matching(); hk.match\_l[u] = partner or -1.
\begin{lstlisting}
HopcroftKarp hk(nl, nr); hk.add_edge(u, v);
int m = hk.max_matching();  // Konig: min vertex cover = m; max independent set = nl + nr - m
\end{lstlisting}
\end{libentry}
\begin{libentry}{Bipartite minimum vertex cover / maximum independent set (Konig)}{Ctrl+Alt+U, V}{cpkonig}{above solve()}{O(V E) (augmenting paths)}{library/flow/min\_vertex\_cover.cpp}
konig(nl, nr, adj): adj[left] = right neighbours. Returns \{cover\_left, cover\_right\}: a minimum vertex cover, whose size equals the maximum matching; the complement is a maximum independent set.
\begin{lstlisting}
auto [cl, cr] = konig(nl, nr, adj);   // vertices in the cover (bool per side)
\end{lstlisting}
\end{libentry}
\begin{libentry}{Global minimum cut (Stoer-Wagner)}{Ctrl+Alt+U, Shift+A}{cpstoerwagner}{above solve()}{O(n\textasciicircum{}3)}{library/flow/stoer\_wagner.cpp}
stoer\_wagner(w): the minimum total weight of edges whose removal disconnects the graph, w = symmetric n x n weight matrix (n \textgreater{}= 2). No source or sink needed.
\begin{lstlisting}
ll cut = stoer_wagner(w);
\end{lstlisting}
\end{libentry}
\begin{libentry}{Gomory-Hu tree (all-pairs minimum cut)}{Ctrl+Alt+U, A}{cpgomoryhu}{above solve()}{n - 1 max flows (Dinic)}{library/flow/gomory\_hu.cpp}
gomory\_hu(n, edges): a tree par[v], w[v] (v \textgreater{}= 1) on 0..n-1 for an undirected capacitated graph; the min cut between any u and v equals the smallest w on the tree path (Gusfield\textquotesingle{}s algorithm, n - 1 max flows).
\begin{lstlisting}
auto [par, w] = gomory_hu(n, edges);   // edges: {u, v, capacity}
\end{lstlisting}
\end{libentry}
\begin{libentry}{Hungarian algorithm (min-cost assignment)}{Ctrl+Alt+U, Shift+H}{cphungarian}{above solve()}{O(n\textasciicircum{}2 m)}{library/flow/hungarian.cpp}
hungarian(a): a is n x m with n \textless{}= m; assigns each row a distinct column minimizing the total cost. Returns \{cost, col\_of\_row\}. For maximum weight, negate the costs.
\begin{lstlisting}
auto [cost, col] = hungarian(a);   // row i -> column col[i]
\end{lstlisting}
\end{libentry}
\begin{libentry}{Max flow (Dinic) + min cut}{Ctrl+Alt+Y, Y}{cpdinic}{above solve()}{O(V\textasciicircum{}2 E) worst case, O(E sqrt V) on unit-capacity graphs}{library/flow/dinic.cpp}
Dinic g(n); g.add\_edge(u, v, cap); ll f = g.max\_flow(s, t); after it g.min\_cut\_side(v) = v on the source side of a min cut.
\begin{lstlisting}
Dinic g(n); int e = g.add_edge(u, v, cap);
ll f = g.max_flow(s, t); ll flow_on_e = g.flow_on(e);
\end{lstlisting}
\end{libentry}
\begin{libentry}{Max flow with lower bounds}{Ctrl+Alt+Y, Shift+G}{cpflowlb}{above solve()}{one or two Dinic runs}{library/flow/lower\_bounds.cpp}
LowerBoundFlow lb(n); lb.add\_edge(u, v, low, high); lb.max\_flow(s, t) = max s-t flow meeting every bound, -1 if none. Needs cpdinic.
\begin{lstlisting}
LowerBoundFlow lb(n); lb.add_edge(u, v, 1, 3);
ll f = lb.max_flow(s, t);  // -1: the bounds cannot be met
\end{lstlisting}
\end{libentry}
\begin{libentry}{Maximum matching in a general graph (Edmonds\textquotesingle{} blossom)}{Ctrl+Alt+U, Shift+G}{cpblossom}{above solve()}{O(n\textasciicircum{}3)}{library/flow/blossom.cpp}
blossom\_matching(adj): match[v] = partner or -1, for any undirected graph (odd cycles allowed). Matching size = number of v with match[v] != -1, divided by 2.
\begin{lstlisting}
vector<int> mate = blossom_matching(adj);
\end{lstlisting}
\end{libentry}
\begin{libentry}{Min-cost max-flow}{Ctrl+Alt+Y, Shift+Y}{cpmcmf}{above solve()}{O(flow * V * E) worst case, far less in practice}{library/flow/mcmf.cpp}
MCMF g(n); g.add\_edge(u, v, cap, cost); auto [flow, cost] = g.run(s, t, limit). Shortest paths by SPFA, so negative costs are fine.
\begin{lstlisting}
MCMF g(n); g.add_edge(u, v, 1, w);
auto [flow, cost] = g.run(s, t);  // cheapest among maximum flows (or among flows of size limit)
\end{lstlisting}
\end{libentry}
\section{Misc}
\begin{libentry}{De Bruijn sequence}{Ctrl+Alt+U, 9}{cpdebruijn}{above solve()}{O(k\textasciicircum{}n)}{library/misc/de\_bruijn.cpp}
de\_bruijn(k, n): a cyclic sequence of length k\textasciicircum{}n over \{0..k-1\} in which every length-n word appears exactly once (read cyclically). Built from Lyndon words (FKM algorithm). Shortest string containing all words: the sequence + its first n-1 symbols.
\begin{lstlisting}
vector<int> seq = de_bruijn(2, 3);   // 0 0 0 1 0 1 1 1
\end{lstlisting}
\end{libentry}
\begin{libentry}{Interactive problem helpers}{Ctrl+Alt+T, Shift+I}{cpinteractive}{inside solve()}{--}{library/misc/interactive.cpp}
ask(i, j) prints "? i j", flushes, reads the reply (exits on -1); answer(x) prints "! x".
\begin{lstlisting}
int r = ask(i, j);
answer(x);
\end{lstlisting}
\end{libentry}
\begin{libentry}{Random generator}{Ctrl+Alt+T, Shift+R}{cprng}{above solve()}{O(1)}{library/misc/random.cpp}
rnd(l, r) uniform in [l, r]; shuffle(v.begin(), v.end(), rng). Seeded by time.
\begin{lstlisting}
ll x = rnd(1, 100);
shuffle(v.begin(), v.end(), rng);
\end{lstlisting}
\end{libentry}
\begin{libentry}{Random hashing of sets and multisets (XOR / sum hashing)}{Ctrl+Alt+U, X}{cpxorhash}{above solve()}{O(1) per label}{library/misc/xor\_hash.cpp}
RandHash h; h(x) = a random 64-bit label for value x. Multisets are equal (whp) iff the sums of labels match; "every value appears an even number of times" iff the XOR of labels is 0. With prefix sums / XORs this answers range questions in O(1).
\begin{lstlisting}
RandHash h; vector<uint64_t> px(n + 1); for (int i = 0; i < n; i++) px[i + 1] = px[i] ^ h(a[i]);
bool all_even = (px[r] ^ px[l]) == 0;   // range [l, r)
\end{lstlisting}
\end{libentry}
\begin{libentry}{Simulated annealing skeleton}{Ctrl+Alt+U, Z}{cpanneal}{above solve()}{(seconds) x (cost of energy + mutate)}{library/misc/annealing.cpp}
anneal(s, energy, mutate, seconds, t0, t1): minimize energy(state) by random local moves; worse moves are accepted with probability exp(-delta / T), T cooling from t0 to t1 geometrically. Returns \{best state, best energy\}. For heuristic / output-only problems.
\begin{lstlisting}
auto [best, e] = anneal(start, [&](const State& s) { return cost(s); }, [&](State s, mt19937_64& rng) { /* small change */ return s; }, 1.8, 10.0, 0.01);
\end{lstlisting}
\end{libentry}
